% Transcription — fonds Alexandre Grothendieck, shelfmark 10, batch 9 (pages 161-180).
% Established from the facsimile archives/batches/10/batch-09.pdf (a mirror of
% https://grothendieck.umontpellier.fr/10.pdf), per the skill
% transcribe-grothendieck.
%
% Pass: Opus 5 (claude-opus-5), 2026-09-12 — first pass, unchecked against the
% pages by a human.
%
% Twenty leaves, all of them on the subject of the carton, all in his own hand,
% and the whole lot is the continuation without a break of the run begun page
% 156 and broken off mid-sentence at the foot of page 160: an abelian class C of
% k-modules, the complete linearly topologised algebra U it reconstructs, and
% the permitted U-modules.
%
% Pages 161-167 carry the supplementary structures: duality as an involutive
% anti-automorphism of U; the tensor product as a diagonal U → U ⊗̂ U; the
% axioms of associativity, commutativity, augmentation and compatibility with
% duality; and page 167's brace, which collects the whole apparatus in four
% articles — a topological algebra, a coassociative diagonal, an augmentation, an
% involution — with P(v ⊗ 1)Δ = P(1 ⊗ v)Δ = 1.
% Pages 168-179 pass to the group: G the invertible s of U with Δs = s ⊗ s and
% ε(s) = 1, then U*, then the primitive elements and their Lie algebra; A = U'
% as an algebra of functions on G; the equivalence of a K-category, a K-bigebra U
% and a commutative hyperalgebra A, with G total in U; and the reduction, through
% an associative algebra Ω and its completion Ω̂ for σ(Ω,A), of the permitted
% modules to ordinary representations. Page 179 is the schema that closes the
% run, K-catégorie ⟺ K-bigèbre U ⟺ hyperalgèbre commutative A, with the
% pseudo-algebraic and the algebraic group hanging below it.
% Page 180 starts afresh, on the free group Z(I) generated by the indecomposable
% classes, ordered by Krull-Schmidt, and on what a tensor product, a privileged
% object and a duality add to it.
%
% Three material facts govern the whole batch and are stated here rather than
% page by page. First, this run is drafting and not copying out: the formulas,
% the underlined statements and the displayed diagrams carry, and a great deal of
% the connective prose does not — pages 161 to 167 in particular are written at
% speed and repeatedly cancelled, and the apparatus is dense in consequence.
% Second, several blocks are cancelled whole, by long diagonals or by a
% hand-drawn loop — the lower two thirds of page 162, the head of page 163 inside
% a hand-drawn box, the middle of page 165, the openings of pages 173 and 175 —
% and they give up their symbols and nothing of the prose between them. Third,
% notes of his run slantwise up the left edge of pages 164 and 168 in a hand
% faster still; they are given as far as they go and no further.
%
%   161  duality as a homomorphism U° → U, i.e. an involutive anti-automorphism
%        x ↦ x̌; compatibility with the augmentation, ε(x̌) = ε(x); and, for a
%        continuous φ : V → U, φ̄(V)' = φ̄(V').
%   162  the tensor product as a bifunctor T over k^U × k^U lying above ⊗_k;
%        U ⊗̂_k U; U_V = U/J_V; the lower two thirds cancelled.
%   163  the indecomposables V_i ⊗ V_j and their W_ijk; W_1 ⊗ W_2 with its
%        « naturelle » structure; the homomorphism U → U ⊗̂ U underlined.
%   164  U/J ⊗_k U/J ⊂ L_k(V) ⊗_k L_k(W) = L_k(W ⊗_k V); the U-multiplications
%        on V ⊗ V as Σ u_i ⊗ v_i; the compatibilities for U → U ⊗̂ U.
%   165  the supplementary axioms: associativity of Δ, commutativity through
%        V ⊗ W → W ⊗ V, the augmentation, and (V ⊗ W)' ≃ V' ⊗ W'; the square of
%        Δ against ∨ ⊗ ∨, i.e. ∨ an isomorphism of coalgebras.
%   166  the combined compatibility of ε, ∨, Δ through V' ⊗ V → k, written
%        P(v ⊗ 1)Δ = 1 and P(1 ⊗ v)Δ = 1; then U, V two « K-algèbres » and the
%        permitted functor φ̄ : k^V → k^U compatible with ⊗ iff φ commutes to Δ.
%   167  the brace collecting the four data on U, and P as the multiplication
%        U ⊗̂ U → U; « Par dualité », the same on an algebra A.
%   168  3) the group G of invertible s with Δs = s ⊗ s and ε(s) = 1;
%        ε(s) = ε(s)², s ∈ G ⇒ ε(s) = 1, and š = s^{-1}; A = U'.
%   169  N.B. the elements of A are the coefficients u ↦ ⟨u·a, b⟩ of the
%        finite-dimensional representations; 1) the group U* of invertibles;
%        3) the primitive elements, Δ(A) = A ⊗ 1 + 1 ⊗ B, and the Lie algebra.
%   170  A as an algebra of functions on G, compatible with the diagonal, the
%        augmentation and ∨; « G est total dans U »; the weak topology;
%        A = k[x]/x^p.
%   171  if G is total in U, k^U is the class of G-modules whose coefficients lie
%        in A; « Renversons le problème »: start from G and A_0 and recover U,
%        with A_0 ⊂ A and G → G(U).
%   172  the identification of G with an algebra A_0 of functions satisfying the
%        preceding conditions; N.B. Ω an associative algebra with 1, A the
%        coefficients of its finite-dimensional representations, the continuity
%        of the multiplication for σ(Ω ⊗ Ω, A ⊗ A) and the homomorphism A → A ⊗ A.
%   173  the finite-dimensional Ω-modules are the Ω̂-modules, Ω̂ the completion of
%        Ω for σ(Ω,A); ⟨u ⊗ v, f⟩ = ⟨uv, f⟩ and f(uv) ∈ A_0 ⊗ A_0.
%   174  the class of Ω-modules whose f_{x,x'} lie in A_0, isomorphic to the
%        A_0-comodules and to the Ω̂-modules for σ(Ω,A_0).
%   175  the coefficients of a representation in V ⊗ W; G and A_0 again.
%   176  CQFD; « Reprenons un groupe G »: k(G), the coefficients, f(uv) in
%        A_0 ⊗ A_0, and the stability of A_0 under multiplication and under ∨.
%   177  1 ∈ A_0 for the trivial representation; the structure G' and the
%        bijectivity of G → G'(A_0).
%   178  V = Σ V_i and the coefficients f_{e_i,e_j}; N.B. the algebra generated
%        by the g_ij and the inverse of the determinant.
%   179  the schema: K-catégorie ⟺ K-bigèbre U ⟺ hyperalgèbre commutative A,
%        with « Groupe pseudo-algébrique » and « Groupe algébrique » below.
%   180  I the set of indecomposable classes, Z(I) the free group it generates,
%        ordered by Krull-Schmidt, I^+ the isomorphism classes; a privileged
%        element e when k carries one, an involution i ↦ ǐ with ě = e, and
%        n(i) = dim M_i extended by linearity.

\documentclass[11pt,a4paper]{article}
% Le préambule commun à toutes les transcriptions du fonds Grothendieck.
%
% Il tient en une idée : **l'incertitude se compose, elle ne se résout pas.**
% Une transcription qui lisse un mot douteux a détruit la seule chose pour
% laquelle on la faisait — un lecteur doit pouvoir savoir, à chaque endroit,
% ce qui était sur le feuillet et ce qui a été ajouté. Les six macros qui
% suivent sont donc l'appareil critique, et non de la décoration.
%
% Les mêmes commandes sont reconnues par `scripts/render.mjs`, qui produit la
% vue de lecture du site. Une macro ajoutée ici doit l'être aussi là-bas,
% faute de quoi le rendu échouera bruyamment — ce qui est voulu : un rendu
% silencieusement faux est pire qu'un rendu absent.

\ProvidesPackage{grothendieck}[2026/08/08 Transcriptions du fonds Grothendieck]

\usepackage{amsmath,amssymb,amsthm}
\usepackage{mathrsfs}
\usepackage{stmaryrd}
% Les diagrammes commutatifs sont partout dans ce fonds : c'est la langue
% même de la géométrie algébrique de Grothendieck, et une transcription qui
% les rendrait en prose ne transcrirait rien.
\usepackage{tikz-cd}
% Français par défaut — c'est la langue des feuillets — mais l'anglais est
% chargé aussi : la traduction et le résumé sont des documents anglais, et la
% typographie française (espaces avant les deux-points) y serait une faute.
% Ces documents basculent par \selectlanguage{english} en tête de corps.
\usepackage[english,french]{babel}
\usepackage{fontspec}
\usepackage{geometry}
\geometry{a4paper,margin=25mm,marginparwidth=28mm}
\usepackage{marginnote}
\usepackage{xcolor}
% ulem, not soul. `soul`'s \st and \ul are built on a character-by-character
% scan that XeLaTeX's font machinery breaks: the document fails with an
% undefined control sequence rather than degrading. ulem does the same two jobs
% and is Unicode-engine safe. `normalem` stops it hijacking \emph, which the
% transcriptions use for Grothendieck's own emphasis.
\usepackage[normalem]{ulem}
\usepackage{hyperref}
\hypersetup{colorlinks=true,linkcolor=black,urlcolor=black,pdfborder={0 0 0}}

% --- Métadonnées du lot -----------------------------------------------------
% Reprises telles quelles par le script de rendu, qui les affiche en tête de
% la vue de lecture. Elles fixent ce que le fichier prétend couvrir : sans
% elles, une transcription flotte sans point d'attache dans un fonds de seize
% mille feuillets.

\newcommand{\folder}[1]{\gdef\grothFolder{#1}}
\newcommand{\batch}[1]{\gdef\grothBatch{#1}}
\newcommand{\leaves}[2]{\gdef\grothFirst{#1}\gdef\grothLast{#2}}
\newcommand{\foldertitle}[1]{\gdef\grothFoldertitle{#1}}
\gdef\grothFolder{?}\gdef\grothBatch{?}\gdef\grothFirst{?}\gdef\grothLast{?}\gdef\grothFoldertitle{}

% --- Le résumé --------------------------------------------------------------
%
% Il ouvre la lecture modernisée, sous le titre « Résumé », et il est composé
% en petit. Ce n'est pas un ornement : c'est le seul passage du document qui
% ne soit pas des mathématiques, et un lecteur doit pouvoir le reconnaître
% avant de l'avoir lu — puis le sauter s'il connaît déjà le sujet, ou s'y
% arrêter s'il ne le connaît pas. Le filet qui le suit marque la reprise du
% corps.

\newenvironment{resume}
  {\small\setlength{\parskip}{0.45em}}
  {\par\medskip\hrule\medskip}

% --- Mots-clés ---------------------------------------------------------------
%
% En anglais, en clôture du résumé de la lecture modernisée : le vocabulaire
% moderne sous lequel chercher ce que les feuillets construisent. Le manifeste
% du site (`npm run manifest`) les extrait pour étiqueter la cote entière —
% cette ligne est la source unique des tags, il n'y a pas de fichier à côté.
\newcommand{\keywords}[1]{\par\smallskip\noindent
  {\footnotesize\textsc{Keywords} --- \textit{#1}}\par}

% --- L'appareil critique ----------------------------------------------------

% Le numéro de feuillet, en marge. C'est l'ancre qui permet de régler un
% désaccord sur une formule en regardant *un* feuillet plutôt qu'une chemise
% de six cents. Le site s'en sert aussi pour tourner les pages du fac-similé
% quand on fait défiler la transcription.
\newcommand{\leaf}[1]{%
  \marginnote{\footnotesize\color{gray!70}\textsf{#1}}%
  \label{leaf:#1}\ignorespaces}

% Illisible : on ne devine pas. Un mot inventé qui se lit comme les autres est
% le pire résultat possible.
\newcommand{\ill}{\textcolor{red!60!black}{[\,\ldots\,]}}

% Lecture douteuse : le mot est donné, et signalé comme incertain.
\newcommand{\uncertain}[1]{\uline{#1}}

% Ajout de l'éditeur — une lettre manquante, un mot sous-entendu.
\newcommand{\add}[1]{\textcolor{blue!50!black}{[#1]}}

% Biffé par Grothendieck lui-même. Ce qu'il a rayé fait partie du document :
% une rature dit souvent où la pensée a changé de direction.
\newcommand{\struck}[1]{\sout{#1}}

% Note du transcripteur, sur l'état matériel du feuillet.
\newcommand{\note}[1]{\par\smallskip{\small\color{gray!60!black}\textit{[#1]}}\par\smallskip}

% Note marginale de Grothendieck, distincte du corps du texte.
\newcommand{\marginal}[1]{\par\smallskip\noindent\hspace{1em}%
  {\small\color{gray!60!black}#1}\par\smallskip}

% --- Titre ------------------------------------------------------------------

\AtBeginDocument{%
  \begin{center}
    {\large\bfseries Fonds Alexandre Grothendieck --- cote n\textsuperscript{o}~\grothFolder}\\[2pt]
    {\itshape\grothFoldertitle}\\[4pt]
    {\small feuillets \grothFirst--\grothLast\ (lot \grothBatch)}\\[2pt]
    {\footnotesize\color{gray!60!black}Transcription établie d'après le fac-similé numérisé
     par l'Université de Montpellier.\\
     Les lectures incertaines sont soulignées, les illisibles notées [\,\ldots\,],
     les ajouts entre crochets.}
  \end{center}
  \vspace{1em}}

\endinput


\folder{10}
\batch{9}
\pages{161}{180}
\dating{[à partir de 1958]}
\watermark{Édition de démonstration}
\foldertitle{Catégories tensorielles : notes manuscrites (s.d.), tapuscrits (s.d.)}

\begin{document}

\section*{Structures supplémentaires (suite) : dualité, produit tensoriel, axiomes}

\note{Les pages 161 à 167 continuent sans coupure la rédaction commencée page
156 et interrompue en pleine phrase au bas de la page 160. Elles sont écrites
d'un trait rapide et abondamment raturées : les formules, les énoncés soulignés
et les diagrammes portent, une grande partie de la prose de liaison ne porte
pas.}

\page{161}

modules dans les $\underline{\mathcal{U}}$-modules, i.e.\ un
\struck{isomorphisme} homomorphisme
$\underline{\mathcal{U}}^{\circ} \to \underline{\mathcal{U}}$, i.e.\ une
\uncertain{application} $\sim\,: \underline{\mathcal{U}} \to
\underline{\mathcal{U}}$ qui est \struck{\ill{}} \ill{} et anti-multiplicative ;
$\check{F}\check{F} = I$ signifie que $\sigma^2 = 1$, donc c'est un \ill{}
\uncertain{anti-automorphisme} involutif, noté $x \to \check{x}$.

\note{Une ligne entière est biffée entre « qui est » et « et
anti-multiplicative » ; elle ne se lit pas.}

Réciproquement, un tel anti-automorphisme définit \struck{\ill{}} un foncteur
$k^{\underline{\mathcal{U}}} \to k^{\underline{\mathcal{U}}}$. Les hypothèses
ci-dessus, qui \ill{} \ill{}. \add{Si l'on s'est donné sur
$\underline{\mathcal{U}}$} une augmentation $\varepsilon$ \ill{} soit
\struck{\ill{}} compatible \ill{}, i.e.
\[ \varepsilon(x) = \varepsilon(\check{x}) , \]
signifie que \struck{\ill{} des $\underline{\mathcal{U}}$-modules} \ill{} les
$\underline{\mathcal{U}}$-modules triviaux. En \uncertain{fait} \ill{}

\ill{} un homomorphisme continu $\varphi : \underline{V} \to
\underline{\mathcal{U}}$, \ill{} compatible avec les \ill{}, \struck{\ill{}}
i.e.
\[ \varphi(\check{x}) = \check{\overline{\varphi(x)}} , \]
signifie que \uncertain{si} \struck{$\overline{\varphi}(V)'$} pour tout
$V \in k^{\underline{\mathcal{U}}}$, \ill{},
\[ \overline{\varphi}(V)' = \overline{\varphi}(V') , \]
i.e.\ le foncteur permis est compatible avec le passage au dual.

\page{162}

\ill{} donner un \emph{bifoncteur} $T(V,W)$ tel que le diagramme

\begin{tikzcd}[column sep=large]
k^{\underline{\mathcal{U}}} \times k^{\underline{\mathcal{U}}}
  \arrow[r, "T"] \arrow[d, "I^{\underline{\mathcal{U}}} \otimes I^{\underline{\mathcal{U}}}"'] &
  k^{\underline{\mathcal{U}}} \arrow[d, "I^{\underline{\mathcal{U}}}"] \\
k^{k} \times k^{k} \arrow[r, "T"'] & k^{k}
\end{tikzcd}

\[ T(V,W) = V \otimes_{k} W \]

soit commutatif. Le \ill{} pour $V, W \in k^{\underline{\mathcal{U}}}$, \ill{}
sur $V \otimes_{k} W$ une structure de $\underline{\mathcal{U}}$-module \ill{},
et telle que \ill{} de $\underline{\mathcal{U}}$-homomorphismes.

\note{Les deux tiers inférieurs du feuillet sont enfermés dans une boucle tracée
à la main et traversés de longues diagonales ; ils ne rendent que les symboles
isolés donnés ici.}

\struck{\ill{}} Soit \ill{}, \ill{}, i.e.
\[ \underline{\mathcal{U}} \mathbin{\hat{\otimes}_{k}} \underline{\mathcal{U}} \]
\ill{} structure de $\underline{\mathcal{U}}$-module ; soit alors
$V \in k^{\underline{\mathcal{U}}}$, $\underline{\mathcal{U}} \to
\underline{\mathcal{U}}(V)$, et \ill{} de $\underline{\mathcal{U}}$ \ill{}
\[ \underline{\mathcal{U}}_{V} = \underline{\mathcal{U}}/J_{V} . \]
\ill{} les $\underline{\mathcal{U}}_{V}$-modules $W$ \ill{}, $W \otimes W$
\ill{},
\[ \underline{\mathcal{U}} \otimes_{k} \underline{\mathcal{U}}(V) \]
\ill{} ; $\underline{\mathcal{U}}_{V}$ étant un anneau \uncertain{artinien}, il
\ill{} de $\underline{\mathcal{U}}_{V}$-modules \ill{}

\page{163}

\note{Le haut du feuillet est enfermé dans un cadre tracé à la main et traversé
de deux longues diagonales.}

\struck{indécomposables distincts, soit $V_i$, $V_j$ deux \ill{}, $V_i \otimes
V_j \to$ un $\underline{\mathcal{U}}$-module, donc la décomposition en une somme
finie de $\underline{\mathcal{U}}$-modules \struck{\ill{}} indécomposables
$W_{ijk}$, l'un des $W_{ijk}$ obtenus est \ill{} ; soit $J'_{V}$ \ill{}
canonique de $\underline{\mathcal{U}}$ dans $\Sigma\, W_{ijk}$}

Montrons qu'il existe un $W \in k^{\underline{\mathcal{U}}}$ \ill{} $V$, $i$,
$j$, $k$. \ill{} $W_1$ et $W_2$ sont $\in k^{\underline{\mathcal{U}}_{V}}$,
alors $W_1 \otimes W_2$, muni de sa structure « naturelle » de
$\underline{\mathcal{U}}$-module, soit un \ill{} $\supset J_{V} W$. En effet, il
suffit de vérifier pour $W_1$ et $W_2$ indécomposables (grâce à
Krull--Schmidt) \ill{} $V$ \ill{} et $V_i$ indécomposables distincts ;
\[ W = \Sigma\, V_i \otimes V_j . \]
Cela dit, \ill{}

\struck{\ill{}}

Ceci est un cas particulier de \uncertain{2} \add{définitions} suivantes. Soit
le $\underline{\mathcal{U}}$-module permis ; \ill{} un module permis sur
$\underline{\mathcal{U}} \mathbin{\hat{\otimes}} \underline{\mathcal{U}}$)
définit un $\underline{\mathcal{U}}$-module permis de façon fonctorielle, \ill{}
\uncertain{un homomorphisme $\underline{\mathcal{U}} \to
\underline{\mathcal{U}} \mathbin{\hat{\otimes}} \underline{\mathcal{U}}$}. Mais,
\ill{}

\page{164}

\ill{} on voit que 2 données \add{de $\underline{\mathcal{U}}$-modules}
\struck{\ill{}} structures sur les $V \otimes_{k} W$
($V, W \in k^{\underline{\mathcal{U}}}$) \ill{}, i.e.\ il \ill{} deux \ill{}
distinctes \ill{} une $\underline{\mathcal{U}} \to \underline{\mathcal{U}}
\mathbin{\hat{\otimes}} \underline{\mathcal{U}}$ \ill{} pour tous \ill{}. Mais,
\struck{\ill{}}, \ill{}, on \ill{} $J$, \ill{} et \ill{}, et \ill{} définirait
$\underline{\mathcal{U}} \to \underline{\mathcal{U}}/J \otimes_{k}
\underline{\mathcal{U}}/J$. ( \ill{} $J$ \ill{} l'\ill{} dans \ill{} d'un
$V \in k^{\underline{\mathcal{U}}}$, donc $\underline{\mathcal{U}}/J =
\underline{\mathcal{U}}_{V} \subset L_{k}(V)$,
\[ \underline{\mathcal{U}}/J \otimes_{k} \underline{\mathcal{U}}/J
   \subset L_{k}(V) \otimes_{k} L_{k}(W) = L_{k}(W \otimes_{k} V) . \]
\struck{\ill{}} \ill{} $\underline{\mathcal{U}} \to L_{k}(W \otimes V)$, on
dirait qu'on \ill{} pour qu'on fasse la \uncertain{répartition} \ill{}

Les $\underline{\mathcal{U}}$-\uncertain{multiplications} sur $V \otimes V$
($V \in k^{\underline{\mathcal{U}}}$) sont de 2 formes
$\Sigma\, u_i \otimes v_i$, \ill{}, où $u_i$, $v_i$ sont des
$\underline{\mathcal{U}}$-\uncertain{multiplications} de $V$.

Alors on obtient $\underline{\mathcal{U}} \to \underline{\mathcal{U}}/J
\otimes_{k} \underline{\mathcal{U}}/J$. On voit \ill{} et les compatibilités
nécessaires pour que $\underline{\mathcal{U}} \to \underline{\mathcal{U}}
\mathbin{\hat{\otimes}} \underline{\mathcal{U}}$ soit défini, et on vérifie que
c'est bon. \ill{} la loi \ill{} chercher.

\marginal{Le \struck{$k^{\underline{\mathcal{U}}} \to$} $k^{\underline{\mathcal{U}}}
\to$ \ill{}, cette \ill{}, et \ill{} vérifiée \uncertain{d'elle-même} \ill{}}

\page{165}

\uncertain{Axiomes supplémentaires.} \uncertain{Associativité} de
l'homomorphisme $\Delta$ : \uncertain{accompagnée d'un axiome} d'associativité
pour les produits tensoriels.

\uncertain{Commutativité de $\Delta$.} L'homomorphisme
$V \otimes_{k} W \to W \otimes_{k} V$ doit être compatible avec les
$\underline{\mathcal{U}}$-structures. Cela ne \ill{} de bifoncteurs tensoriels
\uncertain{monoïdaux} et commutatifs, i.e.\ la \ill{}

\ill{} la compatibilité \add{avec l'augmentation}

\note{Une insertion enfermée dans une boucle tracée à la main occupe ici la
largeur du feuillet ; elle ne rend que la formule qui suit.}

$\Delta(u) = 1 \otimes u + u \otimes 1 + {}$ \ill{}

\struck{\ill{}} \ill{} $V$ \ill{} de $\underline{\mathcal{U}}$-modules
\struck{\ill{} $V$} \ill{} pour $V, W \in k^{\underline{\mathcal{U}}}$,
\[ (V \otimes_{k} W)' \simeq V' \otimes_{k} W' \]
\ill{}, soit \ill{}

\begin{tikzcd}[column sep=large]
\underline{\mathcal{U}} \arrow[r, "\Delta"] \arrow[d, "\vee"'] &
  \underline{\mathcal{U}} \otimes \underline{\mathcal{U}} \arrow[d, "\vee \otimes \vee"] \\
\underline{\mathcal{U}} \arrow[r, "\Delta"'] &
  \underline{\mathcal{U}} \otimes \underline{\mathcal{U}}
\end{tikzcd}

\note{Un bloc entièrement raturé sépare le diagramme de la ligne qui suit ; il
ne se lit pas.}

i.e.\ $\Delta \vee = (\vee \otimes \vee)\Delta$, i.e.\ $\vee$ est un
isomorphisme de co-algèbres.

\page{166}

\uncertain{Compatibilité combinée de $\varepsilon$, $\vee$, $\Delta$.} Ça veut
que pour $V \in k^{\underline{\mathcal{U}}}$, l'homomorphisme naturel
$V' \otimes V \to k$ soit un $\underline{\mathcal{U}}$-homomorphisme, quand $k$
est muni \add{de sa structure} de $\underline{\mathcal{U}}$-module trivial (si
on \ill{} pour supposer $\Delta$ symétrique, il faut évidemment \uncertain{2}
\ill{} pour $V \otimes V' \to k$). Ceci s'exprime par
\[ P(v \otimes 1)\Delta = 1 \quad\text{et}\quad P(1 \otimes v)\Delta = 1 \]
(\ill{} $\Delta$ \ill{} symétrique).

\uncertain{Si tous ces} axiomes sont vérifiés, alors \ill{} des
$\underline{\mathcal{U}}$-modules permis aux permis.

\struck{\ill{}}

Soient $\underline{\mathcal{U}}$, $\underline{V}$ des « $K$-algèbres » avec un
homomorphisme $\Delta$, qui \add{donne} $\underline{V} \to
\underline{\mathcal{U}}$, i.e.\ $\overline{\varphi}$ un foncteur permis
$k^{\underline{V}} \to k^{\underline{\mathcal{U}}}$.

\uncertain{Pour que $\overline{\varphi}$ soit compatible avec les produits
tensoriels} (i.e.\ si $V, W \in k^{\underline{\mathcal{U}}}$, \ill{} sur
$V \otimes W$ la structure de $\underline{V}$-module canonique est la \ill{}
\struck{\ill{}} déduite de $\varphi(V)$ et $\varphi(W)$), il faut et il suffit
que $\varphi$ soit compatible avec $\Delta$.

\page{167}

Si \struck{toutes} les données et tous les axiomes se ramènent réunis sur
$\underline{\mathcal{U}}$, \struck{\ill{}} \ill{} :

\begin{itemize}
\item $\underline{\mathcal{U}}$ algèbre associative topologique quelconque
  \emph{\ill{}} avec suffisamment de représentations de dimension finie ;
\item $\underline{\mathcal{U}}$ muni d'une
  \struck{structure de co-algèbre \ill{}}
  \uncertain{structure de co-algèbre associative}
  $\underline{\mathcal{U}} \to \underline{\mathcal{U}}
  \mathbin{\hat{\otimes}} \underline{\mathcal{U}}$,
  \struck{\ill{} commutative} \ill{} ;
  ces deux structures compatibles, c'est-à-dire un homomorphisme continu
  d'algèbres ;
\item $\underline{\mathcal{U}}$ muni d'une augmentation $\varepsilon$, qui est
  un homomorphisme d'algèbres et un homomorphisme de co-algèbres ;
\item $\underline{\mathcal{U}}$ muni d'une involution $\vee$, qui est un
  anti-\uncertain{automorphisme} (d'algèbres, involutif) et un automorphisme de
  co-algèbres ; $\varepsilon$ invariante sous $\vee$ :
  $\varepsilon(\check{x}) = \varepsilon(x)$.
\end{itemize}

On a
\[ P(v \otimes 1)\Delta = P(1 \otimes v)\Delta = 1 \]
($P$ désigne la multiplication $\underline{\mathcal{U}}
\mathbin{\hat{\otimes}} \underline{\mathcal{U}} \to \underline{\mathcal{U}}$ qui
correspond au produit, \struck{\ill{}} \add{il est bien} \struck{défini ; il
faut que \ill{} pour tout $V \in k^{\underline{\mathcal{U}}}$, \ill{}
$J_{W}\,\underline{\mathcal{U}} \subset$ \ill{}} défini \ill{} (ce n'est
\ill{})).

\uncertain{Par dualité}, ceci correspond aux données \ill{} sur une algèbre
$A$ : \ill{} ; \ill{} de topologie, et \ill{} la structure d'algèbre \ill{}, et
\ill{} des co-algèbres \ill{} $A$ est commutative.

\section*{Le groupe $G$ et l'algèbre $A$}

\page{168}

3) Le groupe $G$ des éléments inversibles de $\underline{\mathcal{U}}$ tels que
\uncertain{$\Delta s = s \otimes s$}, $\varepsilon(s) = 1$. \struck{\ill{}}
\ill{} correspond aux \ill{} des foncteurs $I^{\underline{\mathcal{U}}}$ tels
que si $A$ et $B$ \struck{sont} correspondent à
$V, W \in k^{\underline{\mathcal{U}}}$, alors $A \otimes B$ correspond à
$V \otimes W$. \add{Ils} correspondent aussi \ill{} homomorphismes
\struck{dans} $\underline{\mathcal{U}}$ continus, compatibles avec $\Delta$,
donc \ill{} de l'espace $\underline{\mathcal{U}}(\mathbf{Z})$ \ill{}

\ill{} $\Delta s = s \otimes s$ \ill{}
\[ \varepsilon(s) = \varepsilon(s)^{2} , \]
donc $\varepsilon(s)$ égal à $0$ ou $1$, mais \struck{pas} $0$, car $s$ est
inversible, d'où
\uncertain{$s \in G \Rightarrow \varepsilon(s) = 1$}.
\ill{} $P(v \otimes 1)\Delta = 1$, puis $\check{s}\,s = 1$, d'où
\uncertain{$\check{s} = s^{-1}$}.

\ill{} \uncertain{Conclusion} : si \ill{}, une 2\textsuperscript{e} réduction
\ill{} $\vee : \underline{\mathcal{U}} \to \underline{\mathcal{U}}$, continue,
$P(v \otimes 1)\Delta = 1$. Cette $\vee$ \ill{},
$\varepsilon(\check{x}) = \varepsilon(x)$, \ill{} $s^{\vee} = s^{-1}$, \ill{}
$f \to \langle s, f \rangle$ \ill{} $G$ \ill{}
\[ A = \underline{\mathcal{U}}' \]
\ill{}

\marginal{$G$ \ill{} dans le \ill{} des \ill{} \uncertain{finie} \ill{} la
représentation de base \ill{} \struck{\ill{}} \ill{} $G$ dans $V$ \ill{} de la
représentation \ill{} $G$ dans $V$ \ill{}}

\page{169}

N.B. Les éléments de $A$ \struck{(} i.e.\ les formes linéaires continues sur
$\underline{\mathcal{U}}$) sont exactement les formes du type
\[ u \to \langle u \cdot a,\; b \rangle , \]
où $a, b' \in V \in k^{\underline{\mathcal{U}}}$, i.e.\ les coefficients des
représentations linéaires continues de dimension finie de
$\underline{\mathcal{U}}$. \struck{\ill{}}

\struck{Éléments}

\ill{} \ill{} d'éléments dans $\underline{\mathcal{U}}$ :

1) Le groupe $\underline{\mathcal{U}}^{*}$ \struck{\ill{}} des éléments
inversibles de $\underline{\mathcal{U}}$. Correspond aux \ill{} automorphismes
des foncteurs $I^{\underline{\mathcal{U}}}$, i.e.\ les \struck{\ill{}} \ill{}
pour tout $V \in k^{\underline{\mathcal{U}}}$ d'un \struck{automorphisme} \ill{}
$u_{V}$ de $V$, avec commutation \ill{} pour les
$\underline{\mathcal{U}}$-homomorphismes. \uncertain{\ill{}}

2) \struck{\ill{}} Les éléments primitifs.

3) Le \struck{groupe des éléments} \ill{} Il s'agit des endomorphismes du
foncteur additif $I^{\underline{\mathcal{U}}}$ tels que si $A$ et $B$ sont
\ill{} qui correspondent à $V$ et $W$, \struck{\ill{}}
\[ A \otimes 1 + 1 \otimes B \]
pour $V \otimes W$. Il \ill{} algèbre de Lie \ill{}, par exemple les \ill{} de
Lie \ill{} (\ill{} car $p \ne 0$).

\page{170}

topologique de $\underline{\mathcal{U}}$) dans $k$ (et $=$ \ill{} $k$). Ainsi
\ill{} une \struck{\ill{}} homomorphisme \ill{} d'algèbres
\struck{\ill{} de $A$} dans l'algèbre des fonctions de $G$ à valeurs dans $k$,
transportant $1$ en $\mathbf{1}$, compatible d'ailleurs avec la multiplication
diagonale et l'augmentation, et \ill{} aussi avec $\vee$. Dire que
$G \to$ \ill{} signifie précisément que \uncertain{$G$ est total dans
$\underline{\mathcal{U}}$}.

N.B. \ill{} $G$ est \ill{}, \ill{} les conditions de $A$ \ill{}, signifie \ill{}
l'adhérence \ill{} avec la topologie « faible » \ill{} pour \ill{} $G$ \ill{}.
\ill{} aussi que le \ill{} homomorphisme \ill{} $A$ \ill{} par
$\underline{\mathcal{U}}$, \struck{\ill{}} \ill{} indépendamment \ill{} que $A$
\ill{}

\ill{} (si $x, x' \in V$, $x'_i \in V'$ \ill{}) \ill{} de $V$ contenant \ill{},
\ill{} $V'$ \ill{}
\[ x' \otimes x' \in V' \otimes V' ; \]
\struck{$p \ne 0$} \ill{}
\[ \underline{A} = k[x]/x^{p} \]
\struck{\ill{}} \ill{} $x_i^{p}$ \ill{} l'homomorphisme d'algèbres \ill{}

\page{171}

\struck{\ill{}} i.e.\ $A$ est une algèbre de fonctions sur $G$.

Si $G$ est total dans $\underline{\mathcal{U}}$, alors la classe
$k^{\underline{\mathcal{U}}}$ s'identifie, avec \struck{toutes} ses structures,
à la classe des \struck{\ill{}} $G$-modules \add{tels que} les \ill{}
coefficients sont \struck{\ill{}} \ill{} $\in A$.

Ainsi, une $\underline{K}$-\uncertain{bigèbre} $\underline{\mathcal{U}}$
\ill{} un \struck{\ill{}} groupe $G$, et une \struck{\ill{}} algèbre $A_{0}$ de
fonctions sur $G$ à valeurs dans $k$, séparant les points de $G$, \ill{} stable
\ill{} $\vee$, et telle que
\[ f(xy) \in A_{0} \otimes A_{0} . \]
Alors, $\underline{\mathcal{U}}$ \ill{} définit \ill{} $G$ \ill{} $A_{0}$ est
\struck{\ill{}} \ill{} dans $\underline{\mathcal{U}}$].

\uncertain{Renversons le problème} : \ill{} $G$ avec une algèbre $A_{0}$
vérifiant les conditions précédentes. \ill{} $G$-modules \ill{} $A_{0}$ \ill{}
des formes linéaires \ill{} \ill{} les catégories des $G$-modules \ill{}. Il lui
correspond \ill{} $G \to \underline{\mathcal{U}}$, $A_{0} \subset A$ \ill{}
$\to \underline{\mathcal{U}} \to$ \ill{} une bijection \ill{}
$G(\underline{\mathcal{U}})$ (i.e.\ \ill{} inversibles \ill{} telles que
$\Delta(s) = s \otimes s$). Pour \ill{} par $G$, \ill{} $A_{0}$ \ill{} de $G$,
il \ill{} tel que $G \to G(\underline{\mathcal{U}})$, pour montrer que les
\ill{} \ill{} de $A_{0}$ il suffit de prouver que $A \to A_{0}$

\page{172}

est bijectif. Ceci \ill{} \ill{} \uncertain{plus bas}.

Ainsi, une identification du groupe $G$ avec une algèbre $A_{0}$ de fonctions de
$G$ dans $k$, satisfaisant aux conditions précédentes, et \uncertain{certain}
\struck{\ill{}} $K$-\ill{} $\underline{\mathcal{U}}$.

N.B. Soit $\underline{\mathcal{U}}$ une \ill{} \ill{} de type fini \ill{}, telle
que les $\underline{\mathcal{U}}/J$ \ill{}, $A$ \ill{} la co-algèbre des \ill{}
\struck{\ill{}} de \ill{} co-algèbres de \ill{}). Les
$\underline{\mathcal{U}}$-modules \ill{} des conditions \uncertain{antérieures}
\ill{} linéaires sur $k$.

Soit $\Omega$ une \struck{\ill{}} algèbre associative avec 1, considérons la
classe des $\Omega$-modules \ill{} de dimension finie sur $k$. Soit $A$
l'ensemble des fonctions sur $\Omega$ \ill{} coefficients des représentations
linéaires \ill{} de dimension finie de $\Omega$. La multiplication de $\Omega$
est continue pour $\sigma(\Omega \otimes \Omega, A \otimes A)$ \ill{} et
$\sigma(\Omega, A)$, i.e.\ \ill{} un homomorphisme
\[ A \to A \otimes A , \]
\ill{}, \ill{} ( \ill{} $\Omega = L_{k}(V)$, $V$ de dimension finie sur $k$ ).
Alors, les $\Omega$-modules \ill{} de

\page{173}

dimension finie sont exactement les \uncertain{co-algèbres} \struck{\ill{}}
\ill{} de dimension \ill{} sur $k$. Donc c'est identique : la classe des
$\hat{\Omega}$-modules permis, où $\hat{\Omega}$ est le complété de $\Omega$
pour $\sigma(\Omega, A)$.

\note{Trois lignes entièrement biffées suivent ; elles ne rendent que les mots
« Plus généralement », « avec 1 » et « $A_{0}$ ».}

\struck{Plus généralement \ill{} avec 1, \ill{} $A_{0}$ \ill{}}

\ill{} de $\Omega$ \ill{} les \ill{} \struck{application} \ill{}
\[ \langle u \otimes v,\; f \rangle = \langle uv,\; f \rangle \]
\ill{} \struck{linéaire (qui existe \ill{})} \ill{} application « \ill{} »
\ill{} telles que $\Omega' \otimes \Omega'$ ( \ill{} pour \ill{}
$f(uv) \in A \otimes A$ ; \ill{} si $f(uv) =$ \ill{} alors \ill{}, $f_i$,
$g_i(w) \rangle$, donc \ill{} \ill{} i.e.\ \ill{} Cela dit,
\struck{\ill{}} de $\Omega'$ \ill{} \ill{} il \ill{} que
$f(uv) \in A_{0} \otimes A_{0}$. \ill{} une \ill{} co-algèbre de $A$ ;

\page{174}

\struck{\ill{} $A_{0}$ \ill{} des \ill{}, i.e.\ \ill{}}

Soit $k$ \ill{} classe des $\Omega$-modules \ill{} \add{tels que} pour
$x \in V$, $x' \in V'$, $f_{x,x'}$ soit dans $A_{0}$. Cette classe \ill{} est
isomorphe à la classe des $A_{0}$-\uncertain{comodules} ( \ill{} \ill{}
commutative, \ill{} \ill{} la dimension ; le réciproque \ill{}, grâce à la
dimension finie ). C'est \ill{} la classe des $\hat{\Omega}$-modules, où
$\hat{\Omega}$ est le complété \ill{} $\sigma(\Omega, A_{0})$.

\ill{} \uncertain{la dernière} \ill{}. $A_{0}$ définit \ill{} \ill{} des
\uncertain{dimensions finies} de représentations, \ill{} pour \ill{}, \ill{}
\ill{} de dimension ). \ill{} le groupe $G$. \ill{} les multiplications de
\ill{} de $G$. \struck{Alors} Donc \ill{} conditions précédentes. \ill{}
\ill{} des $\Omega$ dans $V$, \ill{} de $\underline{\mathcal{U}}$
\struck{$f = \Sigma f_i \otimes g_i$, donc \ill{} le \ill{} représentation
\ill{}} \ill{} d'ici une \ill{} algèbre

\page{175}

de plus\,\uncertain{(1)}, une \ill{} \ill{} de \ill{}, \struck{\ill{}} de $V$,
\ill{} $W$, \ill{} coefficient d'une représentation de $V \otimes W$.

\struck{\ill{}}

\ill{} $A_{0}$, \ill{} \ill{} : $G$, \ill{} $A$ \ill{} $A_{0}$. Pourvu \ill{}
par le \ill{} de $G$, contenant $A_{0}$, \struck{\ill{}} \ill{} représentation
de dimension finie de $\Omega$ dans \ill{} dans $G$.
\struck{Dire \ill{} $\in G(A_{0})$ \ill{}} \ill{} dans $A_{0}$ (si \ill{})
(les \ill{} parties \ill{} de $V'$, et \ill{} \ill{} représentation
$\Omega'$-\ill{}, $x \to f_{x,e_i}$) \struck{\ill{} $G(A_{0})$ \ill{} $G$
\ill{}}

\page{176}

d'une représentation dans \ill{} \ill{} sont de \ill{} $u * f$ \ill{}
le représentation de \ill{} $\Omega$ de la représentation de \ill{} dans $G$.
CQFD.

Reprenons \ill{} $G$. \struck{\ill{}} \ill{} une \ill{} représentation $A_{0}$
\ill{} $G^{k}$ et
\[ G^{(k)} = k(G) , \]
et \ill{} l'ensemble des coefficients \ill{} représentations de $G$ \ill{} les
fonctions $f(uv)$, pour $f \in A_{0}$, sont dans \ill{}, \ill{} de \ill{} dans
$A_{0}$ \ill{} $G$. \ill{} que $G$ \ill{} \uncertain{signifie que $A_{0}$ est
stable sous la multiplication}, autrement dit, que \uncertain{$A_{0}$ est une
sous-algèbre}, \ill{} \uncertain{$A_{0}$ est stable sous $\vee$} ;

\page{177}

\ill{} représentation \ill{} triviale dans $k$, \struck{\ill{}} \ill{}
\uncertain{$1 \in A_{0}$} \ill{} \struck{$K$-algèbre} \ill{} dont la \ill{} des
coefficients \ill{} $A_{0}$. \ill{} \struck{\ill{}} \ill{} structure \ill{}
$G'$, et \ill{}
\[ G \to G'(A_{0}) \]
\ill{} bijectif, i.e.\ \ill{} les points de \ill{} fonctions de $A_{0}$, i.e.\
\ill{} sépare les points de $G$, et \ill{} fonctions de $A_{0}$ provient de $G$.
\ill{} condition la structure sur $G$ définie \ill{} le \ill{} de $G$, i.e.\ de
$A_{0}$, \ill{} une structure de \ill{} algébrique \ill{} $A_{0}$ soit
\uncertain{engendrée finiment} sur $k$.

\page{178}

\ill{} suffisamment \ill{} $V$ \ill{} $V_i$, alors
\[ V = \Sigma\, V_i . \]
N.B. \struck{\ill{}} \ill{} \struck{\ill{} de $V$,} \ill{} de $V$,
$f_{e_i, e_j}$, \ill{} \ill{} : $\otimes V$ ( \ill{} ). \ill{} \ill{} une partie
) \ill{} \uncertain{algèbre} \ill{} engendrée par les \ill{} \ill{} \ill{}
l'inverse \ill{} du \uncertain{déterminant} \ill{}

\page{179}

\ill{} \ill{} des représentations \ill{} $L(V)$ \ill{} puis que les fonctions
\ill{}, représentation de $G$, \ill{} \ill{}, \ill{} les fonctions,

\[ K\text{-catégorie} \iff K\text{-bigèbre } \underline{\mathcal{U}}
   \iff \text{hyperalgèbre commutative } A \]

\note{Dans le second membre du schéma, « algèbre » est biffé et « bigèbre »
écrit au-dessus ; le troisième membre se poursuit par deux mots dont seul
« plate » se lit.}

\note{Une flèche verticale monte des deux lignes qui suivent vers le second
membre de ce schéma.}

\uncertain{Groupe pseudo-algébrique} (car si $G(\underline{\mathcal{U}})$
\struck{est total} est total dans $\underline{\mathcal{U}}$, on retrouve
$\underline{A}$ dans $A$ \struck{\ill{}}, i.e.\ \ill{} \ill{} des \ill{} \ill{}
$A \to$ \ill{} : 0)

\uncertain{Groupe algébrique} (si de plus $A$ \ill{} une \ill{} \ill{} de type
fini).

\section*{Le groupe ordonné des classes indécomposables}

\page{180}

Considérons à nouveau une classe abélienne $\mathcal{R}$ dont les \ill{} soient
de longueur finie. Soit $I$ l'ensemble des classes de \ill{} indécomposables
dans $\mathcal{R}$, soit $\mathbf{Z}(I)$ le groupe libre engendré par $I$. C'est
un \uncertain{groupe ordonné} d'après Krull--Schmidt, $I^{+}$ ( l'ensemble des
éléments \ill{} ) est canoniquement isomorphe à l'ensemble des classes de \ill{}
$\in \mathcal{R}$ modulo isomorphisme, la somme de $I^{+}$ correspondant :
l'opération de somme directe.

Si $\mathcal{R}$ est une classe de \struck{\ill{}} \add{\ill{}} modules
\struck{\ill{}} \ill{} structures, $\otimes$ \ill{}, et si \ill{} que $k$ soit
muni d'une structure privilégiée, alors $I$ a un élément distingué $e$ ; si dans
$\mathcal{R}$, \ill{} une involution $\vee$ (le dual d'un module muni d'une
structure \ill{} est muni d'une structure \ill{} $k$) alors $I$ est aussi muni
d'une involution $i \to \check{\imath}$, et la compatibilité de l'involution
avec la structure privilégiée de $k$ s'exprime par
\[ \check{e} = e . \]
N.B. \ill{} sur $I$, il y a une \ill{} \ill{} privilégiée
\[ n(i) = \dim M_{i} , \]
qui se prolonge par linéarité à \ill{} $= \mathbf{Z}(I)$,

\note{Une tache d'encre couvre le symbole précédant l'égalité ; le lot s'arrête
là.}

\end{document}
